\documentclass[a4paper,10pt]{article}

% --- Paquetes ---
\usepackage[utf8]{inputenc}
\usepackage[T1]{fontenc}
\usepackage[spanish]{babel}
\usepackage{graphicx} % Aplicamos draft solo a imágenes para mayor velocidad sin perder enlaces ni funcionalidad
\usepackage{booktabs}
\usepackage{amsmath}
\usepackage{amssymb}
\usepackage{hyperref}
\usepackage{geometry}
\usepackage{caption}
\usepackage{subcaption}
\usepackage{longtable}
\usepackage{array}
\usepackage{float}
\usepackage{titlesec}

% --- Paginación automática ---
\newcommand{\sectionbreak}{\clearpage}
% \newcommand{\subsectionbreak}{\clearpage}

% \newcommand{\subsubsectionbreak}{\clearpage}



% --- Estilo de párrafos con salto de línea ---
\titleformat{\paragraph}[hang]
{\normalfont\normalsize\bfseries}{}{0pt}{}
\titlespacing*{\paragraph}
{0pt}{3.25ex plus 1ex minus .2ex}{1em}

\titleformat{\subparagraph}[hang]
{\normalfont\normalsize\bfseries\itshape}{}{0pt}{}
\titlespacing*{\subparagraph}
{0pt}{2.25ex plus 1ex minus .2ex}{1em}

% --- Configuración de Párrafos (Salto de línea en lugar de sangría) ---
\setlength{\parindent}{0pt}
\setlength{\parskip}{\baselineskip}

% --- Configuración ---
\geometry{margin=1.5cm}
\hypersetup{
    colorlinks=true,
    linkcolor=blue,
    filecolor=magenta,      
    urlcolor=cyan,
    pdftitle={Reporte de Resultados y Comparativa TFG},
}

\title{Reporte de Resultados y Comparativa}
\author{Trabajo de Fin de Grado}
\date{30 de Abril de 2026}

\begin{document}

\maketitle

Este reporte técnico presenta un análisis detallado de la evaluación experimental del Trabajo de Fin de Grado (TFG) centrado en el problema de selección de Tag SNPs mediante algoritmos evolutivos multiobjetivo. El documento comienza con los resultados de referencia de \textbf{Moqa et al. (2022)} y prosigue con la comparativa de dos configuraciones experimentales propias (Ejecuciones \textbf{A} y \textbf{B}), permitiendo validar la eficacia de las propuestas metodológicas desarrolladas.

\begin{itemize}
    \item \textbf{Reporte 1 (Moqa 2022)}: resultados del paper de referencia.
    \item \textbf{Reporte 2 (Ejecución A)}: \texttt{20260430T185224}.
    \item \textbf{Reporte 3 (Ejecución B)}: \texttt{20260430T183144}.
\end{itemize}

Para facilitar la comprensión de las comparativas y los indicadores agregados, se detallan a continuación las tres fuentes de datos analizadas y su tratamiento metodológico respecto a los objetivos de maximización (Tolerancia y Hamming):

\begin{itemize}
    \item \textbf{Estudio Moqa 2022 (Referencia):} Se analizan los resultados extraídos del paper original. La principal problemática de este estudio reside en la dificultad de reproducir los resultados obtenidos por los autores. Este asunto se desarollará en su sección correspondiente. 
    
    Uno de estos problemas de reproducibilidad es la ausencia de una descripción detallada sobre la transformación de objetivos en el estudio. Se hipotetiza que emplean una estrategia de inversión basándose en otro estudio en el que se parcialmente se sustenta: \textbf{Ting et al. (2010)}. En este estudio se invierten los objetivos a maximizar como en la Ejecución B, en vez de negarlos como se hace en la Ejecución A.
    \item \textbf{Ejecución A (Negación):} Emplea una transformación lineal mediante la negación de los objetivos ($f' = -f$). Este enfoque preserva la estructura del espacio original y las distancias relativas entre individuos, lo que permite que métricas como el Hipervolumen mantengan una relación directa y proporcional con el espacio biológico. Esta idea metodológica fue propuesta por el tutor de este trabajo.
    \item \textbf{Ejecución B (Inversión):} Utiliza una transformación no lineal mediante la inversión de los valores ($f' = 1/f$). Esta técnica penaliza severamente los valores bajos de precisión, obligando a los algoritmos a priorizar soluciones de alta calidad biológica, aunque introduce una distorsión en la topología del frente de Pareto que dificulta la comparación directa de magnitudes con la Ejecución A.
\end{itemize}


\newpage
{
\setlength{\parskip}{0pt}
\tableofcontents
}

% =========================================================
\section{Reporte 1 — Moqa et al., 2022}
\subsection{Configuración resumida (Moqa)}
En este bloque se sintetiza el estado del conocimiento sobre la configuración técnica del estudio de referencia, distinguiendo entre los datos verificables y aquellos que han sido inferidos o permanecen como incógnitas.


 \paragraph{Datos Confirmados}
 \begin{itemize}
     \item \textbf{Tamaño del haplotipo:} 1032 SNPs (Dataset Hinds et al., 2005).
     \item \textbf{Representación:} Cadena binaria.
     \item \textbf{Tamaño de la población:} 200 individuos.
     \item \textbf{Tamaño de la descendencia:} 200 individuos.
     \item \textbf{Tamaño del vecindario ($T$):} 15 (específico para MOEA/D).
     \item \textbf{Probabilidad de cruce ($p_c$):} 0,7.
     \item \textbf{Probabilidad de mutación ($p_m$):} $1/l$, donde $l$ representa la longitud del vector de decisión.
     \item \textbf{Criterio de finalización:} 500 generaciones.
     \item \textbf{Estrategia de inicialización:} Aleatoria (Random) y Heurística (Greedy).
     \item \textbf{Número de algoritmos:} 4 (NSGA-II, SPEA2, NSGA-III y MOEA/D).
     \item \textbf{Número de ejecuciones:} 5 por cada configuración experimental.
     \item \textbf{Número total de experimentos:} 40.
 \end{itemize}


\subsection{Profundización de la configuración}
\label{sec:marco_repro_moqa}

\subsubsection{Especulaciones e inconsistencias}
\textit{Limitaciones metodológicas y omisiones de reporte que imposibilitan la replicación exacta del estudio.}

A pesar de los resultados reportados, el estudio de Moqa et al. (2022) presenta una serie de ambigüedades técnicas y vacíos de información que se detallan a continuación:

\paragraph{Hallazgos y Especulaciones}
A través del análisis detallado del código fuente y la documentación técnica complementaria, se han podido esclarecer diversos aspectos de la implementación que permanecían opacos en el paper original, aunque muchos de ellos son meras especulaciones. Estas interpretaciones, derivadas de la inspección del código base de algunos estudios similares y la literatura de referencia, permiten realizar una comparativa más justa y fundamentada:

Mediante la ingeniería inversa del código de Ting et al. (2010), se han identificado los siguientes aspectos:

\begin{itemize}
    \item \textbf{Infraestructura de Datos y Preprocesado:} En el estudio de Ting et al. (2010), se utiliza un bloque genómico de 1032 SNPs en el que se realiza un filtrado automático de marcadores monomórficos antes de la optimización; dicho procedimiento se ha replicado fielmente en este trabajo. El número efectivo de variables de decisión queda definido, por tanto, por el subconjunto de marcadores polimórficos tras este proceso.
    \item \textbf{Lógica de la Inicialización Greedy:} Ante la falta de profundidad en la descripción de la inicialización \textit{greedy} de Moqa, se ha optado por implementar la lógica de Ting en este estudio (además de otras estrategias). El análisis del código de Ting et al. (2010) revela que su inicialización greedy construye cromosomas de forma incremental hasta alcanzar una cobertura del 100\% de los pares de haplotipos e incorpora técnicas estocásticas adicionales. Se profundizará en este asunto en secciones posteriores.
    \item \textbf{Gestión de Semillas:} La naturaleza no determinista de las ejecuciones de Ting se atribuye a una inicialización estocástica dependiente del estado del generador, lo que vincula la reproducibilidad al esquema de siembra aleatoria basado en el tiempo del sistema. Se ha implementado esta metododología.
    \item \textbf{Transformación de Objetivos:} La metodología de Ting contempla transformaciones de inversión ($1/f^2$ y $1/f^3$) para convertir los objetivos de maximización en problemas de minimización tratables por el motor evolutivo. Dado que es plausible que Moqa empleara un enfoque similar, se han evaluado dos tipos de transformación de objetivos: mediante negación y mediante inversión.
    \item \textbf{Normalización Proporcional de Objetivos:} Mediante el estudio detallado de la propuesta de Ting et al. (2010), se ha identificado que la distancia de Hamming media ($f_3$) y su varianza ($f_4$) se evalúan como proporciones relativas al número de marcadores seleccionados ($k$), una transformación fundamental que no se declara explícitamente en la formulación matemática original. Esta métrica se ha integrado en el presente estudio para garantizar la comparabilidad con dicha implementación. 
    
    Es altamente probable que se haya empleado esta técnica, ya que en experimentos previos (como el del 25\_04\_2026), en los que no se aplicó dicha normalización, los resultados divergían significativamente de los de Moqa. Además, esta técnica aporta mayor equidad al evaluar la calidad de las soluciones.
    
    Bajo esta metodología, los objetivos transformados ($f'_3$ y $f'_4$) corresponden a:
    \begin{equation}
        f'_3 = \frac{\frac{1}{N_{pair}} \sum_{i=1}^{N_{pair}} D_i}{k} \quad ; \quad f'_4 = \frac{\text{Var}(D)}{k^2}
    \end{equation}
    Donde $D_i$ representa la distancia de Hamming para el par $i$ sobre el subconjunto de $k$ SNPs seleccionados.
\end{itemize}



\paragraph{Inconsistencias y Lagunas Críticas}
Más allá de las interpretaciones técnicas, el estudio de referencia adolece de una serie de omisiones críticas en su reporte que comprometen la integridad de la replicación. Estas lagunas no solo afectan a la configuración paramétrica de los algoritmos, sino que impactan directamente en la interpretación biológica de los resultados presentados:

\begin{itemize}
    \item \textbf{El Problema del Hipervolumen (HV):} Existe una falla de coherencia interna grave: los autores anuncian el HV como su métrica principal en el \textit{Abstract}, pero los resultados no aparecen en ninguna sección del cuerpo del documento. Además, se omite la declaración del punto de referencia (\textit{Reference Point}), lo que imposibilita cualquier comparación futura o validación de la calidad del frente de Pareto.
    \item \textbf{Indeterminación del Protocolo de Normalización:} El artículo indica explícitamente que \guillemotleft todos los objetivos son normalizados\guillemotright, pero omite por completo la fórmula matemática empleada (Z-score, Min-Max, Inversión, etc.). Esta carencia es crítica, pues la normalización no solo es un paso previo necesario para métricas como el Hipervolumen, sino que altera la topología del espacio de búsqueda y el comportamiento de los algoritmos de ordenamiento no dominado.
    \item \textbf{Crítica a la Trazabilidad y Relevancia Biológica:} Aunque el estudio es reproducible numéricamente gracias a la existencia de un archivo \texttt{input.txt}, carece de anclaje genómico real, es decir, de valor científico biológico. El manuscrito omite el número de cromosoma, las coordenadas exactas (ensamblajes GRCh37/38) y los identificadores rsID. Resulta especialmente problemático el reporte de que \textit{\guillemotleft Haplotype patterns are combined together to form one haplotype block\guillemotright} (pág. 9), lo que sugiere la creación de un bloque artificial mediante la mezcla de poblaciones (europeos, africanos y asiáticos). Dado que el Desequilibrio de Ligamiento (LD) es específico de cada población, esta práctica genera una estructura sintética que no existe en el genoma humano real, reduciendo el experimento a un ejercicio de optimización matemática sin validez biológica contrastable.
    \item \textbf{Ambigüedad Paramétrica en MOEA/D:} No se declara la función de agregación utilizada (Weighted Sum, PBI o Tchebycheff). En caso de haber usado PBI, se desconoce el parámetro de penalización ($\theta$), y para NSGA-III no se especifica el número de divisiones ($H$) que determina la densidad de los vectores de referencia.
    Por este motivo, en las ejecuciones A y B se evalúan de forma exhaustiva diferentes funciones de agregación para el algoritmo MOEA/D: 
    \begin{itemize}
        \item \textbf{WS (Weighted Sum):} Realiza una agregación lineal simple de los objetivos.
        \item \textbf{TCHE (Tchebycheff):} Minimiza la distancia de Tchebyshev respecto a un punto de referencia, permitiendo alcanzar soluciones en frentes no convexos.
        \item \textbf{PBI (Penalty-based Boundary Intersection):} Busca un equilibrio entre convergencia y diversidad mediante el ajuste del parámetro de penalización ($\theta$).
    \end{itemize}
\end{itemize}



Esta falta de transparencia refuerza la importancia de la implementación propia desarrollada en este trabajo, la cual garantiza una especificación completa de parámetros y una semántica de métricas rigurosa.

\subsection{Resultados y análisis (Moqa 2022)}
\label{sec:res_analisis_moqa}

\subsubsection{Frentes de Pareto}
En esta sección se presentan los frentes de Pareto reportados por Moqa et al. (2022), contrastando el desempeño de los algoritmos bajo inicialización aleatoria (RI) y heurística (GI).

\begin{figure}[H]
    \centering
    \begin{subfigure}[b]{0.45\textwidth}
        \centering
        \includegraphics[width=\textwidth]{analisis_assets/moqa_2022/figures/moqa_fig_nsga2_random.jpeg}
        \caption{NSGA-II (RI)}
    \end{subfigure}
    \hfill
    \begin{subfigure}[b]{0.45\textwidth}
        \centering
        \includegraphics[width=\textwidth]{analisis_assets/moqa_2022/figures/moqa_fig_nsga2_greedy.jpeg}
        \caption{NSGA-II (GI)}
    \end{subfigure}
    \vskip\baselineskip
    \begin{subfigure}[b]{0.45\textwidth}
        \centering
        \includegraphics[width=\textwidth]{analisis_assets/moqa_2022/figures/moqa_fig_spea2_random.jpeg}
        \caption{SPEA2 (RI)}
    \end{subfigure}
    \hfill
    \begin{subfigure}[b]{0.45\textwidth}
        \centering
        \includegraphics[width=\textwidth]{analisis_assets/moqa_2022/figures/moqa_fig_spea2_greedy.jpeg}
        \caption{SPEA2 (GI)}
    \end{subfigure}
    \caption{Comparativa de Frentes de Pareto: NSGA-II y SPEA2 (Moqa 2022).}
\end{figure}

\begin{figure}[H]
    \centering
    \begin{subfigure}[b]{0.45\textwidth}
        \centering
        \includegraphics[width=\textwidth]{analisis_assets/moqa_2022/figures/moqa_fig_nsga3_random.jpeg}
        \caption{NSGA-III (RI)}
    \end{subfigure}
    \hfill
    \begin{subfigure}[b]{0.45\textwidth}
        \centering
        \includegraphics[width=\textwidth]{analisis_assets/moqa_2022/figures/moqa_fig_nsga3_greedy.jpeg}
        \caption{NSGA-III (GI)}
    \end{subfigure}
    \vskip\baselineskip
    \begin{subfigure}[b]{0.45\textwidth}
        \centering
        \includegraphics[width=\textwidth]{analisis_assets/moqa_2022/figures/moqa_fig_moead_random.jpeg}
        \caption{MOEA/D (RI)}
    \end{subfigure}
    \hfill
    \begin{subfigure}[b]{0.45\textwidth}
        \centering
        \includegraphics[width=\textwidth]{analisis_assets/moqa_2022/figures/moqa_fig_moead_greedy.jpeg}
        \caption{MOEA/D (GI)}
    \end{subfigure}
    \caption{Comparativa de Frentes de Pareto: NSGA-III y MOEA/D (Moqa 2022).}
\end{figure}

\subsubsection{Convergencia progresiva}
\begin{figure}[H]
    \centering
    \begin{subfigure}[b]{0.48\textwidth}
        \centering
        \includegraphics[width=\textwidth]{analisis_assets/moqa_2022/figures/moqa_fig_convergencia_1.jpeg}
        \caption{Evolución por generación de seis métricas de rendimiento en el experimento base (500 generaciones, población = 200). Se ilustra el impacto de la inicialización heurística (GI, líneas discontinuas) frente a la aleatoria (RI, líneas continuas).}
    \end{subfigure}
    \hfill
    \begin{subfigure}[b]{0.48\textwidth}
        \centering
        \includegraphics[width=\textwidth]{analisis_assets/moqa_2022/figures/moqa_fig_convergencia_2.jpeg}
        \caption{Estudio de sensibilidad de parámetros genéticos (300 generaciones, población = 100). Análisis focalizado en la robustez de los algoritmos ante una reducción de recursos computacionales.}
    \end{subfigure}
    \caption{Análisis de las trayectorias de convergencia reportadas por Moqa et al. (2022).}
\end{figure}

\subsubsection{Top 10 por métrica (Moqa)}
A continuación se presentan los valores finales extraídos de las gráficas de convergencia del estudio original. Estos datos representan el rendimiento absoluto alcanzado en la generación 500.

\subparagraph{Top 10: Range ($\uparrow$)}
\begin{table}[H]
    \centering
    \begin{tabular}{cllc}
        \toprule
        Pos & Algoritmo & Inicialización & Valor ($\mu \pm \sigma$) \\
        \midrule
        1 & MOEAD & R & 2.75 $\pm$ 0.3981 \\
        2 & MOEAD & G & 2.60 $\pm$ 0.4080 \\
        3 & NSGA2 & G & 2.50 $\pm$ 0.4522 \\
        4 & NSGA2 & R & 2.45 $\pm$ 0.4359 \\
        5 & NSGA3 & G & 2.45 $\pm$ 0.4390 \\
        6 & SPEA2 & R & 1.80 $\pm$ 0.2499 \\
        7 & SPEA2 & G & 1.75 $\pm$ 0.2128 \\
        8 & NSGA3 & R & 1.10 $\pm$ 0.0923 \\
        \bottomrule
    \end{tabular}
    \caption{Top para Range en el reporte de Moqa (Valores extraídos de gráfica).}
\end{table}

\subparagraph{Top 10: SumMin ($\downarrow$)}
\begin{table}[H]
    \centering
    \begin{tabular}{cllc}
        \toprule
        Pos & Algoritmo & Inicialización & Valor ($\mu \pm \sigma$) \\
        \midrule
        1 & MOEAD & G & 0.29 $\pm$ 0.0435 \\
        2 & NSGA3 & G & 0.32 $\pm$ 0.0443 \\
        3 & MOEAD & R & 0.35 $\pm$ 0.0603 \\
        4 & NSGA2 & G & 0.35 $\pm$ 0.0483 \\
        5 & NSGA2 & R & 0.38 $\pm$ 0.0542 \\
        6 & SPEA2 & G & 0.41 $\pm$ 0.0354 \\
        7 & SPEA2 & R & 0.46 $\pm$ 0.0506 \\
        8 & NSGA3 & R & 0.60 $\pm$ 0.0249 \\
        \bottomrule
    \end{tabular}
    \caption{Top para SumMin en el reporte de Moqa (Valores extraídos de gráfica).}
\end{table}

\subparagraph{Top 10: MinSum ($\downarrow$)}
\begin{table}[H]
    \centering
    \begin{tabular}{cllc}
        \toprule
        Pos & Algoritmo & Inicialización & Valor ($\mu \pm \sigma$) \\
        \midrule
        1 & MOEAD & G & 0.23 $\pm$ 0.0179 \\
        2 & MOEAD & R & 0.27 $\pm$ 0.0238 \\
        3 & NSGA2 & G & 0.32 $\pm$ 0.0229 \\
        4 & NSGA2 & R & 0.34 $\pm$ 0.0296 \\
        5 & SPEA2 & G & 0.37 $\pm$ 0.0252 \\
        6 & SPEA2 & R & 0.39 $\pm$ 0.0404 \\
        7 & NSGA3 & G & 0.54 $\pm$ 0.0390 \\
        8 & NSGA3 & R & 0.57 $\pm$ 0.0243 \\
        \bottomrule
    \end{tabular}
    \caption{Top para MinSum en el reporte de Moqa (Valores extraídos de gráfica).}
\end{table}

\subparagraph{Top 10: Max. Tolerance Rate ($\uparrow$)}
\begin{table}[H]
    \centering
    \begin{tabular}{cllc}
        \toprule
        Pos & Algoritmo & Inicialización & Valor ($\mu \pm \sigma$) \\
        \midrule
        1 & NSGA3 & G & 0.30 $\pm$ 0.0204 \\
        2 & NSGA3 & R & 0.29 $\pm$ 0.0160 \\
        3 & MOEAD & R & 0.275 $\pm$ 0.0149 \\
        4 & MOEAD & G & 0.27 $\pm$ 0.0173 \\
        5 & NSGA2 & G & 0.25 $\pm$ 0.0170 \\
        6 & NSGA2 & R & 0.24 $\pm$ 0.0146 \\
        7 & SPEA2 & G & 0.23 $\pm$ 0.0108 \\
        8 & SPEA2 & R & 0.23 $\pm$ 0.0113 \\
        \bottomrule
    \end{tabular}
    \caption{Top para Max. Tolerance Rate en el reporte de Moqa (Valores extraídos de gráfica).}
\end{table}

\subparagraph{Top 10: Avg. Tolerance Rate ($\uparrow$)}
\begin{table}[H]
    \centering
    \begin{tabular}{cllc}
        \toprule
        Pos & Algoritmo & Inicialización & Valor ($\mu \pm \sigma$) \\
        \midrule
        1 & MOEAD & R & 0.26 $\pm$ 0.0099 \\
        2 & MOEAD & G & 0.26 $\pm$ 0.0130 \\
        3 & NSGA3 & R & 0.21 $\pm$ 0.0120 \\
        4 & NSGA3 & G & 0.21 $\pm$ 0.0075 \\
        5 & NSGA2 & G & 0.20 $\pm$ 0.0075 \\
        6 & NSGA2 & R & 0.20 $\pm$ 0.0060 \\
        7 & SPEA2 & G & 0.19 $\pm$ 0.0054 \\
        8 & SPEA2 & R & 0.19 $\pm$ 0.0056 \\
        \bottomrule
    \end{tabular}
    \caption{Top para Avg. Tolerance Rate en el reporte de Moqa (Valores extraídos de gráfica).}
\end{table}

\subparagraph{Top 10: Avg. Hamming ($\uparrow$)}
\begin{table}[H]
    \centering
    \begin{tabular}{cllc}
        \toprule
        Pos & Algoritmo & Inicialización & Valor ($\mu \pm \sigma$) \\
        \midrule
        1 & SPEA2 & G & 0.94 $\pm$ 0.0098 \\
        2 & SPEA2 & R & 0.91 $\pm$ 0.0094 \\
        3 & NSGA2 & R & 0.90 $\pm$ 0.0191 \\
        4 & NSGA2 & G & 0.87 $\pm$ 0.0195 \\
        5 & NSGA3 & G & 0.80 $\pm$ 0.0177 \\
        6 & NSGA3 & R & 0.78 $\pm$ 0.0191 \\
        7 & MOEAD & G & 0.76 $\pm$ 0.0326 \\
        8 & MOEAD & R & 0.75 $\pm$ 0.0309 \\
        \bottomrule
    \end{tabular}
    \caption{Top para Avg. Hamming en el reporte de Moqa (Valores extraídos de gráfica).}
\end{table}

\subsubsection{Ranking Promedio (Average Rank)}
El análisis del ranking global mediante el método \textit{Average Rank} (promedio de las posiciones obtenidas en las 6 métricas reportadas) permite identificar las configuraciones más equilibradas del estudio original con una métrica estadísticamente más robusta.

\begin{table}[H]
    \centering
    \begin{tabular}{lllc}
        \toprule
        Pos & Algo & Init & Avg. Rank (Friedman) \\
        \midrule
        1 & MOEAD & G \\
        2 & MOEAD & R \\
        3 & NSGA3 & G \\
        4 & NSGA2 & G \\
        5 & NSGA2 & R \\
        6 & SPEA2 & G \\
        7 & NSGA3 & R \\
        8 & SPEA2 & R \\
        \bottomrule
    \end{tabular}
    \caption{Ranking Global (Average Rank) para el reporte de Moqa et al. (2022). Un \textbf{valor menor} indica un desempeño superior.}
\end{table}

\subsubsection{Conclusiones del Reporte Moqa}
 \begin{itemize}
     \item \textbf{Mejores métodos:} El análisis mediante \textit{Average Rank} sitúa a \textbf{MOEA/D+G} en la \textbf{primera posición} del ranking global de Moqa, seguido por \textbf{MOEA/D+R}. Esto respalda, en términos ordinales, la robustez de los algoritmos basados en descomposición en el estudio original.
     \item \textbf{Mejores inicializaciones:} La inicialización \textbf{Greedy (G)} tiende a ocupar posiciones superiores frente a variantes aleatorias, especialmente en combinaciones con MOEA/D y NSGA-III, según el ranking promedio.
     \item \textbf{Otros apuntes a destacar:} A nivel de métricas individuales, \textbf{SPEA2} destaca en el ranking de diversidad (\textit{Avg. Hamming Distance}) dentro del reporte original, lo cual motiva su consideración cuando la variabilidad de Tag SNPs es prioritaria.
 \end{itemize}

\subsubsection{Comparativa con las Conclusiones de los Autores}
Al contrastar el análisis del TFG con las declaraciones literales de Moqa et al. (2022), se procede a la validación de sus conclusiones fundamentales:

\begin{itemize}
    \item \textbf{Superioridad de los algoritmos de muchos objetivos:} Los autores sostienen que MOEA/D y NSGA-III ofrecen un rendimiento superior a NSGA-II y SPEA2. Se valida íntegramente para MOEA/D (posiciones 1 y 2 del ranking). En el caso de NSGA-III, la superioridad se valida únicamente bajo inicialización heurística (G, posición 3), mientras que bajo inicialización aleatoria su desempeño decae drásticamente.
    \item \textbf{Liderazgo de MOEA/D:} Afirman que MOEA/D supera al resto en la mayoría de los casos. Se valida de forma absoluta, ya que tanto MOEA/D+G como MOEA/D+R encabezan el \textit{Average Rank} global, demostrando ser la arquitectura más robusta para este problema.
    \item \textbf{Impacto de la inicialización Greedy:} Sostienen que la inicialización heurística mejora los resultados para NSGA-III, SPEA2 y MOEA/D. Se valida categóricamente; el ranking muestra una mejora sistemática en los rangos promedio al emplear G frente a R para todos los algoritmos citados.
    \item \textbf{Validación de la diversidad en SPEA2:} A pesar de no ser el más equilibrado globalmente, se mantiene la validación de la afirmación específica de los autores sobre la superioridad de SPEA2 en la preservación de la distancia de Hamming, una vez analizados los valores finales de las gráficas.
\end{itemize}

